\section{课程定位：什么是深度强化学习}

CS224R 的第一讲从最根本的问题出发：什么是深度强化学习（Deep Reinforcement Learning），它与传统的机器学习有什么不同，以及为什么值得学习。

\begin{importantbox}{深度强化学习的定义}
深度强化学习研究的是\textbf{序贯决策问题}（Sequential Decision-Making Problems）：系统需要基于信息流不断地"观察——行动——观察——行动"。课程同时研究这类问题的\textbf{解法}，包括 imitation learning、online/offline RL、model-free/model-based RL、multi-task/meta RL 以及 RL for LLMs 和 RL for robots。所有方法都强调能扩展到深度神经网络。
\end{importantbox}

\subsection{与监督学习的核心区别}

在监督学习中，我们拥有标注数据 $\{(x_i, y_i)\}$，目标是学习一个函数 $f(x) \approx y$。数据点之间是独立同分布（i.i.d.）的，并且我们被\textbf{直接告知}正确的输出。

在强化学习中，情况完全不同：

\begin{itemize}
    \item \textbf{目标}：学习行为策略 $\pi(a \mid s)$，将状态映射到动作。
    \item \textbf{反馈}：不是直接告诉答案，而是通过\textbf{间接反馈}（如奖励信号）从经验中学习。
    \item \textbf{数据非 i.i.d.}：智能体的动作 $a$ 会影响未来的观察 $s'$，因此数据分布取决于当前策略。
\end{itemize}

\begin{knowledgebox}{行为的多种形态}
深度强化学习中的"行为"可以包括：机器人的运动控制、聊天机器人的对话、游戏的策略、自动驾驶的决策、Web agent 的操作等。这些场景的共同特点是：模型的输出（动作）会对环境产生后果，进而影响未来的输入。
\end{knowledgebox}

\subsection{课程范围与目标}

Chelsea Finn 明确指出，这门课不可能覆盖深度强化学习的所有内容。课程聚焦于：
\begin{itemize}
    \item 深度 RL 方法背后的\textbf{核心概念}，以便学生能理解更高级的方法。
    \item 算法的\textbf{实现}，使学生能够自己编程实现。
    \item 以\textbf{机器人控制和语言模型}为主要示例，但技术具有更广泛的适用性。
\end{itemize}

\subsection{本章小结}

深度强化学习是关于序贯决策的学科。与监督学习相比，RL 面临的独特挑战包括：反馈是间接的、数据是非 i.i.d. 的、动作有后果。课程的核心目标是让学生能够理解和实现现有和新兴的深度 RL 方法。

